\part{Les archives historiques de la Faculté de médecine de Montpellier :
	histoire, constitution et organisation d'un patrimoine universitaire}

\chapter{Enseigner et gouverner la médecine à Montpellier :
	institutions et production documentaire du Moyen Âge à nos jours}
\chapitrecourt{Enseigner et gouverner la médecine}

\section{Montpellier, une ville favorable à l'émergence d'un enseignement médical}
\sectioncourte{L'émergence d'un enseignement médical}

\subsection{La formation de la ville médiévale}

À la différence de nombreuses villes du Midi méditerranéen, Montpellier ne s'est développée ni sur le site d'une cité antique ni autour d'un ancien siège épiscopal. Mentionné pour la première fois en 985, le lieu n'est documenté de manière suivie qu'à partir du dernier tiers du XI\textsuperscript{e} siècle, lorsque s'affirme la seigneurie des Guilhem dans un Bas-Languedoc où s'enchevêtrent notamment les pouvoirs du comte de Melgueil et de l'évêque de Maguelone\footcite{DuhamelAmado1991}\footcite{Katsura1992}. La \emph{villa}, encore périphérique parmi les possessions du lignage autour de l'an mil, devient progressivement sa résidence principale et le centre de son pouvoir\footcite{Vidal1986Donation}\footcite{DuhamelAmado1991}.

Une part importante de cette histoire est connue par le \emph{Liber instrumentorum memorialis}, cartulaire seigneurial compilé au tournant des XII\textsuperscript{e} et XIII\textsuperscript{e} siècles. Il conserve notamment la copie de l'acte du 26 novembre 985 par lequel Bernard II, comte de Melgueil, et son épouse Sénégonde concèdent à un certain Guilhem des biens situés à Montpellier et à Candillargues\footcite[fol.~29~v°, \emph{Carta donacionis feudi, quod contulit Bernardus comes, cum uxore sua, Guillelmo, domino Montispessulani}, vue~36/225 de la version numérisée]{MontpellierAMAA1}. Donations, ventes, arbitrages, serments et privilèges y ont été réunis et ordonnés afin de conserver les titres de la seigneurie et d'en permettre la mobilisation en cas de contentieux ; sa composition résulte toutefois d'une sélection opérée parmi une documentation antérieure en partie dispersée ou perdue\footcite{Chastang2006Liber}.

Aux XI\textsuperscript{e} et XII\textsuperscript{e} siècles, l'affirmation des Guilhem accompagne une croissance rapide de Montpellier. Plusieurs noyaux d'habitat progressivement réunis forment un espace urbain où se développent les activités artisanales et commerciales, notamment autour de la Condamine et de Notre-Dame-des-Tables\footcite{FabreLochard1988}\footcite{Peyron1979}. Une première enceinte réunit plusieurs de ces secteurs dans les années 1130--1140 ; à partir de 1196, la construction et l'entretien de la Commune Clôture sont confiés à une administration spécialisée, qui fixe durablement le périmètre de la ville médiévale\footcite{FabreLochard1988}\footcite{Chastang2013Montpellier}.

Cette croissance s'accompagne d'une évolution des formes de gouvernement. Les bourgeois prennent au cours du XII\textsuperscript{e} siècle une place croissante dans l'entourage et l'administration seigneuriale\footcite{Vergos2021Guilhem}, tandis que s'affirme l'\emph{universitas}, communauté des habitants reconnue comme personne collective et progressivement dotée de ses propres institutions\footcite{Chastang2013Montpellier}. À la fin du XII\textsuperscript{e} siècle, Montpellier est ainsi devenue l'un des principaux centres urbains du Bas-Languedoc, dont le développement politique et institutionnel accompagne une insertion croissante dans les réseaux d'échanges méditerranéens.

\subsection{Un contexte commercial et économique favorable}

L'essor de Montpellier repose en partie sur son insertion dans les échanges méditerranéens. Située à proximité du littoral sans être directement exposée sur une côte ouverte, la ville communique avec les lagunes et la mer par le Lez et les installations portuaires de Lattes, tout en se trouvant au croisement de voies terrestres reliant les espaces ibériques, provençaux et italiens\footcite{FabreLochard1988}\footcite{Galano2021}.

Cette prospérité ne résulte toutefois pas de la seule géographie. Les Guilhem contrôlent marchés, monnaies et prélèvements sur les échanges, favorisent l'aménagement de quartiers et la réouverture du port de Lattes. Dès 1080, des hommes de Montpellier négocient avec le vicomte et l'archevêque de Narbonne au sujet des leudes et de l'accès au port de Cabrela ; à la fin du XI\textsuperscript{e} siècle, un marché est attesté dans la Condamine. Au siècle suivant, l'essor de l'artisanat, notamment textile, et d'une bourgeoisie marchande modifie l'entourage et les pratiques de gouvernement des seigneurs\footcite{FabreLochard1988}\footcite{Vergos2021Guilhem}.

Benjamin de Tudèle, de passage vers 1165, présente Montpellier comme un \enquote{lieu fort commode pour le trafic}, fréquenté par des marchands venus notamment d'Égypte, du Levant, de Lombardie, de Rome, de France, d'Espagne et d'Angleterre, et souligne la présence des Génois et des Pisans. Son récit atteste l'étendue des relations commerciales autant que l'importance de la communauté juive montpelliéraine, dont il mentionne plusieurs savants\footcite[p.~7, vue~16/512 de la version numérisée]{BenjaminTudele1830}.

Ces réseaux relient aussi la Méditerranée aux centres d'échanges septentrionaux. Au milieu du XIII\textsuperscript{e} siècle, une lettre adressée depuis Paris par Firmin de Posquières à l'épicier montpelliérain Nicolas Vézian mentionne une commande de sucre rosat d'Alexandrie, d'épices, de drogues, d'eaux parfumées et de préparations destinées aux foires de Champagne\footcite{MontpellierAMEE289}\footcite{Bach2020}. Les épiciers montpelliérains participent donc à la redistribution de marchandises orientales entre les ports méditerranéens, Montpellier, Paris et les foires champenoises.

Plusieurs de ces produits entrent dans la composition des remèdes médiévaux. Les échanges approvisionnent les épiciers et apothicaires montpelliérains en épices et en drogues venues du bassin méditerranéen, dont une partie est transformée sur place en préparations médicinales\footcite{Irissou1947}. Sans démontrer un transfert direct de connaissances médicales, cette circulation les met en relation avec des produits, des techniques et des intermédiaires appartenant à des espaces linguistiques et culturels différents.

Les fonds municipaux attestent la persistance de cet horizon oriental. En 1264, Jacques Ier d'Aragon accorde une lettre de rémission à des marchands montpelliérains établis à Alexandrie\footcite{MontpellierAMLouvet4247}. Le \emph{Petit Thalamus}, livre municipal constitué par le consulat à partir du XIII\textsuperscript{e} siècle et réunissant textes normatifs, serments, listes consulaires et chronique urbaine, rapporte le tremblement de terre qui endommage le phare d'Alexandrie en 1303\footcite[fol.~25 ; vue~55/1140 de la version numérisée]{MontpellierAMAA9}\footcite{Chastang2013Montpellier}. En 1366, Urbain V autorise les marchands de Montpellier à envoyer annuellement, pendant six ans, un navire vers Alexandrie, à condition qu'il ne transporte ni armes, ni bois, ni fer, ni autres marchandises stratégiques\footcite{MontpellierAMLouvet2799}. Ces actes n'établissent pas une activité commerciale uniforme, mais la présence durable d'Alexandrie dans l'horizon économique et documentaire montpelliérain.

Marchands, clercs, juristes et voyageurs mettent ainsi la ville en contact avec des langues, des pratiques et des savoirs divers. Cette ouverture n'explique pas seule l'apparition de l'enseignement médical, mais elle en favorise les conditions sociales et matérielles : présence de praticiens, accès aux substances médicinales et insertion dans des réseaux dépassant le Bas-Languedoc.

\subsection{Les premiers enseignements médicaux à Montpellier}

\subsubsection{Une réputation médicale antérieure à l'institution}

L'enseignement médical précède son organisation institutionnelle. Plusieurs témoignages attestent dès le XII\textsuperscript{e} siècle la réputation des praticiens montpelliérains, sans permettre de restituer précisément leurs écoles ou leurs méthodes. Les statuts promulgués par Conrad d'Urach en 1220 ne créent donc pas la médecine à Montpellier : ils organisent collectivement des maîtres et des étudiants déjà présents\footcite{Moulinier2003Montpellier}\footcite{Verger2021}.

Le premier témoignage connu figure dans la \emph{Vita Adalberti II}, composée par Anselme de Havelberg. Adalbert de Sarrebruck, alors jeune clerc, se rend à Montpellier pour y étudier la \emph{physica}, terme désignant ici la médecine. La ville est présentée comme la demeure de celle-ci et ses médecins comme des praticiens autant que des maîtres transmettant leur doctrine\footcite[p.~592--593]{Jaffe1867}\footcite{Touati2018}.

La lettre adressée en 1153 par Bernard de Clairvaux à Héraclius de Montboissier renseigne plutôt la réputation thérapeutique de la ville. Tombé malade sur la route de Rome, l'archevêque de Lyon est conduit à Montpellier, où il dépense pour les médecins \enquote{ce qu'il avait et ce qu'il n'avait pas} (\emph{quod habebat et quod non habebat})\footcite[col.~512B]{BernardClairvauxEpistola307}\footcite[46]{Touati2018}.

Jean de Salisbury fournit, dans la seconde moitié du siècle, une indication plus directement scolaire. Au livre I du \emph{Metalogicon}, il évoque ceux qui, renonçant aux exigences de la philosophie, se rendent \enquote{à Salerne ou à Montpellier} (\emph{Salernum vel ad Montempessulanum profecti}) pour apprendre la médecine, avant de revenir rapidement se prévaloir du titre de praticien\footcite{JeanSalisburyMetalogicon1929}. Son intention est satirique : il dénonce les individus qui dissimulent une formation superficielle sous un vocabulaire savant. Le rapprochement avec Salerne, dont l'école de médecine jouit alors d'une renommée majeure en Occident\footcite{DeDivitiisEtAl2004}, n'en situe pas moins Montpellier parmi les destinations reconnues pour l'apprentissage médical. Ces sources font donc apparaître une progression : renommée thérapeutique, attraction d'étudiants étrangers, puis identification de la ville comme centre d'enseignement.

\subsubsection{Un milieu méditerranéen de circulation des savoirs}

Cette affirmation accompagne le renouvellement de la médecine savante aux XI\textsuperscript{e} et XII\textsuperscript{e} siècles. Les traductions effectuées notamment dans la péninsule Ibérique et en Italie méridionale rendent accessibles en latin des textes grecs et arabes qui renouvellent l'étude d'Hippocrate et de Galien et diffusent les œuvres de Rhazès et d'Avicenne. Elles fournissent aux écoles occidentales un corpus théorique plus étendu, étudié par la lecture et le commentaire des autorités médicales\footcite{Siraisi1990}\footcite{Moulinier2003Montpellier}.

La position commerciale de Montpellier facilite la circulation des hommes et des livres, sans prouver l'utilisation de textes déterminés dans ses écoles. Les sources antérieures au XIII\textsuperscript{e} siècle renseignent mieux la réputation du centre médical que le contenu des leçons : l'environnement méditerranéen doit donc être distingué de la transmission effectivement documentée des œuvres au sein du cursus\footcite{Touati2018}.

Les communautés juives du Bas-Languedoc participent à cet environnement. Narbonne, Lunel, Posquières, Béziers et Montpellier forment un réseau de circulation des lettrés, manuscrits et traductions entre arabe et hébreu\footcite{Sirat1979}\footcite{Iancu2025}. Installé à Lunel au milieu du XII\textsuperscript{e} siècle, Judah ibn Tibbon traduit des œuvres philosophiques et morales de Baḥya ibn Paquda, Juda Halevi et Salomon ibn Gabirol. Son fils Samuel achève à Montpellier, en 1204, sa traduction du \emph{Guide des égarés} de Maïmonide ; Moïse ibn Tibbon y traduit au XIII\textsuperscript{e} siècle de nombreux textes philosophiques, scientifiques et médicaux\footcite{Sirat1979}\footcite{Dumas2019Flux}.

Ces activités attestent l'existence locale d'un milieu capable de recevoir et de transformer des savoirs issus de la tradition arabe, mais non une collaboration institutionnelle entre traducteurs juifs et maîtres chrétiens, ni le passage direct des traductions hébraïques dans l'enseignement latin. Elles doivent être situées à l'échelle d'un espace intellectuel régional favorable aux transferts plutôt que présentées comme une origine directement documentée de l'école médicale\footcite{Touati2018}\footcite{Iancu2025}.

\subsubsection{Des écoles libres à l'\emph{universitas medicorum}}

L'organisation des écoles montpelliéraines avant 1220 demeure mal connue. Salerne fournit un point de comparaison, non la preuve d'une filiation institutionnelle : aux XI\textsuperscript{e} et XII\textsuperscript{e} siècles, son enseignement repose sur des maîtres autonomes, choisis et probablement rémunérés par leurs élèves, sans université unifiée ni cursus imposé par une autorité commune\footcite{VergerSalerne1999}\footcite{Verger2021}.

La formation savante s'appuie progressivement sur l'\emph{Articella}, corpus réunissant notamment l'\emph{Isagoge} de Johannitius, les \emph{Aphorismes} et les \emph{Pronostics} d'Hippocrate, ainsi que les traités sur les urines et le pouls attribués à Théophile et Philaret\footcite{Articella1998}\footcite{Siraisi1990}\footcite{OBoyle2000}. Devenu un fondement de l'enseignement médical occidental, cet ensemble a pu circuler avec les maîtres et les manuscrits entre Salerne et Montpellier, sans que l'on puisse en déduire une organisation ou un programme identiques\footcite{VergerSalerne1999}\footcite{Moulinier2003Montpellier}.

Les témoignages du XII\textsuperscript{e} siècle font néanmoins apparaître plusieurs médecins enseignant avant l'existence d'une autorité corporative. La charte accordée en janvier 1181 par Guilhem VIII donne à cette situation sa première formulation juridique connue\footcite[fol.~96, vue~100/225 de la version numérisée]{MontpellierAMAA1}\footcite{GuilhemVIII1181}. Elle interdit de réserver l'enseignement à un maître et garantit à chacun la liberté d'ouvrir une école : elle ne fonde ni faculté ni cursus, mais protège la pluralité des enseignants contre toute appropriation exclusive\footcite{DarricauLugat1999}\footcite{Moulinier2003Montpellier}.

La charte de 1181 reconnaît ainsi une activité scolaire sans instituer de cursus, d'examen commun ni d'autorité corporative. L'organisation collective de l'enseignement médical relève des textes adoptés à partir de 1220\footcite{Verger2021}\footcite{LavabreBertrand2020}.

\section{La structuration institutionnelle de la Faculté}

\subsection{Les fondements juridiques de l'Université de médecine (1220-1289)}

\subsubsection{Les statuts de Conrad d'Urach : l'organisation juridique de la communauté médicale}

La liberté reconnue en 1181 ne suffisait pas à organiser collectivement maîtres et étudiants. Le 17 août 1220, Conrad d'Urach, cardinal-évêque de Porto et de Sainte-Rufine et légat d'Honorius III dans le Midi, promulgue des statuts destinés à assurer la conservation de l'enseignement médical (\emph{ad conservationem medicinalis studii}). Leur préambule rappelle son rayonnement antérieur en évoquant les médecins montpelliérains actifs \enquote{dans diverses parties du monde} (\emph{in diversis mundi partibus})\footcite[fol.~187--188, vues~186--187/348]{ADHG1126}\footcite{Conrad1220}\footcite{Verger2021}\footcite{LavabreBertrand2020}.

Nul ne peut désormais enseigner publiquement sans avoir été examiné et approuvé par l'évêque de Maguelone assisté de maîtres régents (\emph{nisi prius examinatus fuerit et approbatus}). La possibilité d'ouvrir une école subsiste, mais son exercice devient subordonné à une validation collective. Chaque étudiant doit en outre être placé sous l'autorité d'un maître déterminé (\emph{certi magistri sit addictus regimini}) et les enseignants ne peuvent détourner les élèves de leurs confrères (\emph{nullus magistrorum scolarem alterius scienter alliciat vel sollicitet}). Les premiers organes de gouvernement sont également institués.\footnote{Sur le gouvernement de la Faculté, voir la \cref{sec:gouvernement-faculte}.}

Les statuts ne remplacent donc pas les écoles par un établissement centralisé mais ils leur superposent une procédure d'habilitation, des règles communes et des organes de gouvernement. Ils constituent l'acte d'organisation de l'\emph{universitas medicorum}, progressivement complété au cours du XIII\textsuperscript{e} siècle\footcite{Verger2021}\footcite{Moulinier2003Montpellier}.

Les statuts complémentaires des 14 et 21 janvier 1240, adoptés à la suite d'un différend entre maîtres, maintiennent les privilèges accordés par les légats Conrad et Guy (\emph{exceptis que in privilegiis... continentur}) et précisent l'accès à la maîtrise, les serments et la participation aux assemblées. Ils imposent aussi le prêt d'un ouvrage médical authentique lorsqu'aucun exemplaire ne peut être acquis, organisant l'accès aux textes nécessaires à l'enseignement\footcite{AHFMB1}\footcite{Statuts1240}.

\subsubsection{La confirmation pontificale des statuts par Alexandre IV}

Le 28 février 1258, Alexandre IV confirme depuis Viterbe les dispositions de 1220 : interdiction d'enseigner publiquement sans examen et approbation de l'évêque assisté des régents, et obligation d'observer les autres statuts de Conrad\footcite{AHFMB2}\footcite{AlexandreIV1258}. Le pape les déclare \enquote{valides et fermes} (\emph{rata et firma}) et les place sous la garantie apostolique (\emph{presentis scripti patrocinio communimus}). La bulle intègre donc ces règles au corpus des privilèges pontificaux et accroît leur autorité juridique.

\subsubsection{La création du \emph{studium generale} par Nicolas IV}

Le 26 octobre 1289, la bulle \emph{Quia Sapientia} de Nicolas IV institue à Montpellier un \emph{studium generale} réunissant droit, médecine et arts\footcite{MontpellierAMLouvet619}\footcite{NicolasIV1289}. La médecine entre ainsi dans un établissement pontifical commun sans perdre son autonomie corporative. Le pape justifie cette décision par la réputation d'une ville \enquote{très célèbre et renommée} (\emph{celebris plurimum et famosus}) et particulièrement propre à l'étude (\emph{aptus valde pro studio}).

La bulle maintient l'examen des candidats par l'évêque de Maguelone ou son délégué, assisté des maîtres de la faculté concernée, et confère aux licenciés le droit d'enseigner partout sans nouvel examen (\emph{absque examinatione vel approbatione alia, regendi et docendi ubique plenam et liberam habeant facultatem}). Ce \emph{ius ubique docendi} donne aux grades montpelliérains la portée générale propre aux \emph{studia generalia} reconnus par la papauté\footcite{Verger2021}\footcite{Moulinier2003Montpellier}. \emph{Quia Sapientia} réunit ainsi plusieurs disciplines dans un même cadre pontifical et étend la validité de leurs grades.

\subsection{Le gouvernement de la Faculté et l'exercice de ses privilèges}
\label{sec:gouvernement-faculte}

L'Université de médecine est une corporation dotée de compétences normatives, administratives et juridictionnelles, mais insérée dans un système de tutelles. Les régents gouvernent l'enseignement et les promotions ; le chancelier rend la justice et convoque les assemblées ; le doyen exerce la préséance parmi les maîtres en activité ; les procureurs administrent les biens communs. L'évêque de Maguelone participe à la collation de la licence et à la désignation du chancelier, tandis que papauté, seigneurs puis monarchie française confirment les privilèges, arbitrent les conflits et garantissent leur exécution. Cette distribution des pouvoirs produit les statuts, délibérations, actes de grades, comptes, privilèges, nominations et règlements qui documentent l'ordre institutionnel autant que son fonctionnement effectif.

\subsubsection{Le collège des maîtres et la certification des compétences}

Les maîtres régents sont les docteurs habilités à enseigner publiquement. Ils forment le corps professoral, siègent dans les assemblées et évaluent les candidats. Le grade médiéval sanctionne certes des connaissances mais, surtout, il reconnaît publiquement une compétence, fixe une position dans la corporation et ouvre, au terme du cursus, l'accès au magistère\footcite{Verger2023Examens}. Le baccalauréat marque l'entrée dans l'enseignement sous l'autorité des maîtres ; la licence est l'autorisation juridique d'enseigner accordée par l'autorité ecclésiastique après examen ; la maîtrise ou le doctorat confère solennellement les insignes et prérogatives correspondants\footcite{Verger2023Examens}. La promotion associe donc jugement des pairs et validation extérieure.

Les statuts de 1240 précisent la procédure instituée en 1220. Le candidat doit normalement avoir étudié la médecine pendant au moins trois ans et demi, accompli six mois de pratique hors de Montpellier et participé aux disputes. Un maître ès arts reçu à Paris ou dans un autre \emph{studium} réputé peut être présenté après deux ans et demi d'études médicales\footcite{AHFMB1}\footcite{Statuts1240}.

Le certificat délivré le 2 juin 1260 à Guilhem Rotberti atteste l'application de cette procédure à l'autorisation de pratiquer. Deux maîtres délégués par l'évêque l'interrogent sur les matières pratiques, puis rendent compte de sa compétence, de son aptitude, de sa réputation et de sa fréquentation assidue des maîtres\footcite[fol.~14~v°--15~r°, vues~22--23/194]{ADHG1277}\footcite{Certificat1260}. L'évaluation porte donc moins sur l'accomplissement formel d'une durée d'études que sur des savoirs, une expérience et une \emph{fama}.

Après l'extension géographique de la licence par \emph{Quia Sapientia}, les constitutions de Clément V du 8 septembre 1309 renforcent le contrôle du cursus. La première fixe les ouvrages à étudier et exige cinq années de médecine des maîtres ès arts reçus dans un \emph{studium} réputé, contre six pour les autres candidats, chaque année étant comptée pour huit mois. Elle impose aussi une période de pratique, la participation aux disputes et deux leçons publiques, théorique et pratique. Après le retrait du candidat, les maîtres déclarent sous serment s'il est suffisant ou insuffisant, sans comparaison avec les autres (\emph{absque omni comparacione})\footcite[fol.~81~v°]{AHFMS1}\footcite{ClementVAuteurs1309}\footcite{Verger2023Examens}. Une seconde constitution rend la licence nulle sans l'accord de l'évêque et des deux tiers des maîtres résidant dans le \emph{studium} (\emph{nullius penitus sit momenti})\footcite[fol.~83~v°]{AHFMS1}\footcite{ClementVLicence1309}. Les statuts de 1335 réservent enfin aux maîtres reçus la participation aux actes publics de licence et de doctorat (\emph{soli magistrati})\footcite[fol.~68~v°, vue~76/194]{ADHG1277}\footcite{Statuts1335}.

Les statuts de 1340 codifient l'ensemble du parcours : matricule, baccalauréat, examens public et secret, licence, maîtrise, durées d'études, lectures, pratique, disputes, serments, cérémonies et obligations des régents\footcite{Statuts1340}. La matricule enregistre les dates d'arrivée et de départ pour vérifier la durée effective de la formation. En 1383, des statuts complémentaires règlent l'ordre de présentation des candidats et confirment qu'une licence ou une maîtrise ès arts acquise avant le baccalauréat médical dispense d'une année d'études ; ils encadrent également les repas, les paiements et les présents afin de limiter le coût des promotions et les risques de corruption\footcite{Statuts1383ADH}\footcite{Statuts1383AHFM}.

La multiplication des règles n'empêche pas les conflits. En 1389, le choix du maître chargé de promouvoir un candidat, les présents qui lui sont remis et les dépenses du doctorat donnent encore lieu à un arbitrage de l'évêque de Maguelone\footcite{Compromis1389}. L'organisation universitaire se construit ainsi par l'articulation des normes générales et de leur correction contentieuse. Les statuts de 1554 adaptent enfin ce dispositif au collège des docteurs et règlent exercices, examens, promotions, serments, obligations des régents et répartition des droits\footcite[fol.~160v--164v]{AHFMS1}\footcite[transcription des lettres patentes de Louis XII relatives à la fondation des docteurs stipendiés, p.~107--109]{Astruc1767}.

La collation des grades est donc à la fois une fonction académique, un privilège corporatif et une activité documentaire : statuts pour les règles, matricules et registres pour les cursus, certificats et procès-verbaux pour les examens, actes de réception et lettres de grades pour les promotions. Ces séries donnent à voir la manière dont la Faculté définit, vérifie et certifie la compétence médicale.

\subsubsection{Les offices de la communauté universitaire}

Le gouvernement courant repose sur trois offices complémentaires, qui distinguent progressivement fonctions judiciaires, académiques et financières.

\paragraph{Le chancelier : juridiction et convocation de la communauté}

Les statuts de 1220 instituent un chancelier de l'université des étudiants (\emph{cancellarius Universitatis scolarium}), choisi parmi les maîtres\footcite[fol.~188, vue~187/348]{ADHG1126}\footcite{Conrad1220}. Il rend la justice aux maîtres et aux étudiants dans les causes civiles, les affaires criminelles restant réservées à l'évêque de Maguelone. En 1240, il reçoit seul le pouvoir de convoquer les assemblées générales tenues sous serment et jure devant l'évêque et l'Université de rendre une justice impartiale, d'observer les statuts et de les transmettre à ses successeurs\footcite{AHFMB1}\footcite{Statuts1240}. Juridiction, convocation et conservation de la norme font ainsi de la chancellerie le principal office permanent de la communauté.

En 1309, Clément V subordonne la désignation du chancelier à l'accord de l'évêque et des deux tiers des maîtres résidant dans le \emph{studium}\footcite[fol.~80~r°]{AHFMS1}\footcite{ClementVChancelier1309}. La papauté peut néanmoins réserver l'office : le 23 avril 1323, Jean XXII le confère directement à Bertrand Portal \enquote{avec tous ses droits et dépendances}, en écartant les statuts et coutumes contraires\footcite{JeanXXII1323}. Cette provision exceptionnelle limite ponctuellement l'élection collégiale sans abolir les institutions internes de la Faculté\footcite{Verger1991UniversitesMidi}.

Les actes suivants montrent la fonction en exercice. En 1335, Raymond de Molières convoque les maîtres par cédule et préside l'assemblée qui règle commémoration des défunts, émoluments professoraux et promotions\footcite[fol.~68~v°, vue~76/194]{ADHG1277}\footcite{Statuts1335}. Les statuts de 1340 lui attribuent la convocation de deux assemblées ordinaires annuelles et des réunions extraordinaires, ainsi que la surveillance de l'application des règles\footcite{Statuts1340}. En 1364, le cardinal Jean de Saint-Marc, commissaire d'Urbain V, annule la procédure irrégulière ayant conduit à la désignation de Jean Jacobi avant de lui conférer lui-même l'office ; pour l'avenir, il réaffirme l'élection par les maîtres sous confirmation épiscopale, tout en réservant au Saint-Siège le pouvoir de réformer les statuts\footcite[fol.~8~v°, vues~15--16/194]{ADHG1277}\footcite{JeanSaintMarc1364}. La chancellerie se situe donc au croisement des prérogatives corporatives, épiscopales et pontificales.

\paragraph{Le doyen : ancienneté et préséance académique}

Les statuts de 1340 attribuent le décanat au plus ancien des maîtres enseignant effectivement ; en son absence ou en cas de maladie, le régent suivant par ancienneté le remplace jusqu'à son retour\footcite{Statuts1340}. Le doyen bénéficie de la préséance dans le collège et les cérémonies et peut suppléer le chancelier lorsque celui-ci est absent ou que l'office est vacant.

La décision de Jean de Saint-Marc du 7 octobre 1364 confirme qu'il doit être le maître le plus ancien dans le grade parmi ceux qui enseignent ordinairement dans le \emph{studium} (\emph{antiquior in gradu magister, legens ordinarie in dicto studio predicte Facultatis, sit decanus})\footcite[fol.~8~v°, vues~15--16/194]{ADHG1277}\footcite{JeanSaintMarc1364}. Son serment d'assurer personnellement sa lecture annuelle, sauf empêchement légitime, lie ainsi préséance académique et exercice effectif de la régence.

\paragraph{Les procureurs : biens, comptes et représentation}

Les statuts de 1340 prévoient la désignation annuelle de deux procureurs selon une rotation associant les régents les plus anciens aux plus récemment reçus. Ils sont chargés de conserver les biens communs et de rendre compte de leur administration au terme de leur mandat\footcite{Statuts1340}.

En 1554, les procureurs perçoivent les droits versés lors des examens et promotions et les répartissent entre les docteurs\footcite[fol.~160v--164v]{AHFMS1}\footcite[p.~126--129]{Astruc1767}. Ils peuvent également représenter l'Université et participer à la production de certains actes. Astruc signale ainsi d'anciennes lettres de baccalauréat munies d'un \emph{Sigillum Procuratoris Universitatis Medicorum Studii Monspessulani}, dont il déduit prudemment qu'elles pouvaient être expédiées par les procureurs\footcite[p.~89]{Astruc1767}. Sans suffire à reconstituer la procédure, ce sceau suggère leur participation à la production des actes de grades.

Cédules, procès-verbaux, comptes, actes scellés et correspondances matérialisent ainsi l'exercice quotidien de fonctions distinctes mais coordonnées.

\subsubsection{Défendre et étendre les privilèges de la Faculté}

\paragraph{Des droits garantis par plusieurs autorités}

L'autonomie de l'Université reste relative : les maîtres évaluent les candidats, délibèrent et administrent les biens ; l'évêque intervient dans la licence et la chancellerie ; la papauté confirme les privilèges, arbitre les conflits et peut réserver des offices. Les interventions de Jean XXII, Benoît XII et Urbain V dans les universités méridionales associent soutien aux études, contrôle institutionnel et recrutement de gradués au service de la cour pontificale\footcite{Verger1991UniversitesMidi}.

Ces droits doivent être régulièrement confirmés et rendus opposables. Les bulles de Clément V relatives au chancelier et à la licence font ainsi l'objet, en 1324 et 1326, de copies notariées destinées à préserver les originaux tout en permettant leur production devant plusieurs juridictions\footcite{ConfirmationChancelier1324}\footcite{ConfirmationLicence1326}. Cette pratique, qui sera étudiée dans le chapitre consacré à la constitution des archives, montre déjà que l'exercice d'un privilège dépend de la conservation et de l'authentification de son titre.

La puissance séculière concourt parallèlement à leur application. Après l'acquisition de Montpellier par Philippe VI en 1349, qui fait passer la seigneurie des rois de Majorque\footnote{La seigneurie de Montpellier passe des Guilhem à la maison d'Aragon en 1204 par le mariage de Marie de Montpellier avec Pierre II d'Aragon, puis échoit au XIII\textsuperscript{e} siècle à la branche des rois de Majorque\footcite{Irissou1947}.} au roi de France\footcite{Viard1910Montpellier}, les actes royaux interdisent l'exercice médical sans examen ni licence, ordonnent aux officiers de poursuivre les contrevenants et protègent les universitaires contre certaines impositions. En 1395, Charles VI confirme ainsi les exemptions accordées par Charles V en 1380 ainsi que les lettres de Jacques Ier interdisant la pratique irrégulière de la médecine\footcite{CharlesVIPratique1395AHFM}\footcite{CharlesVIPratique1395ADH}\footcite{CharlesVIExemptions1395}.

Ces immunités provoquent des conflits avec les pouvoirs urbains. En 1352, les consuls dénoncent l'extension abusive de la juridiction du Petit Sceau, chargée de protéger les universitaires. Le roi maintient leurs privilèges mais rappelle que leur gardien ne peut soustraire indûment les habitants aux juridictions ordinaires\footcite{PetitScel1352}. L'Université ne se situe donc pas hors de la ville : ses droits se négocient dans un espace où se superposent juridictions corporative, épiscopale, municipale et royale.

\paragraph{Réglementer les professions de santé}

Le contrôle de la pratique médicale constitue l'un des principaux privilèges universitaires. D'abord centré sur les médecins, il s'étend aux professions qui préparent ou administrent les remèdes. Les statuts de 1340 organisent la surveillance des apothicaires et interdisent la vente de certains médicaments sans conseil médical ; les actes royaux de 1376 et de 1399 soumettent médecins, chirurgiens et apothicaires, chrétiens comme juifs, à des exigences de compétence et d'approbation\footcite{Statuts1340}\footcite{CharlesVIPratique1399}. La Faculté justifie ce contrôle par la défense conjointe de sa réputation et du bien public face aux dangers imputés aux praticiens insuffisamment formés.

\paragraph{Assurer les ressources de l'enseignement}

La viabilité du \emph{studium} suppose aussi des maîtres, des étudiants et des revenus durables\footcite{Verger1991UniversitesMidi}. En 1369, Urbain V fonde le Collège des Douze-Médecins afin d'entretenir douze étudiants en médecine originaires de la ville et du diocèse de Mende\footcite{UrbainVCollege1369}\footcite{UrbainVCollege1369Germain}. Les difficultés de perception suscitent confirmations, protections et procédures judiciaires sous Grégoire XI, Clément VII et Benoît XIII, montrant que la pérennité d'une fondation dépend de l'administration de ses revenus et de la défense de ses titres\footcite{GregoireXICollege1371ADH}\footcite{GregoireXICollege1371Germain}\footcite{ClementVIICollege1392ADH}\footcite{ClementVIICollege1392Germain}\footcite{BenoitXIIICollege1395ADH}\footcite{BenoitXIIICollege1395Germain}.

À partir de la fin du XV\textsuperscript{e} siècle, la monarchie agit plus directement sur l'organisation matérielle de l'enseignement. Les lettres patentes de Charles VIII de mai 1496 confirment et amplifient les privilèges de l'Université ; un mandement du 9 janvier 1497 en ordonne la vérification et l'enregistrement\footcite{AHFMB25}. Le 29 août 1498, Louis XII institue quatre docteurs stipendiés sur les finances du Languedoc afin d'assurer des leçons permanentes. Le roi désigne les premiers titulaires ; les vacances suivantes doivent être pourvues par l'évêque de Maguelone et les régents après examen\footcite[fol.~75~v°--77~v°]{AHFMS1}\footcite{AHFMB26}\footcite[p~107--109]{Astruc1767}. Cette mesure distingue progressivement les professeurs rémunérés pour un service régulier des autres docteurs appartenant au collège\footcite{LavabreBertrand2021Esprit}.

Henri II confirme les statuts de 1554, puis la spécialisation des régences se poursuit avec l'anatomie et la démonstration des simples en 1593, la chirurgie et la pharmacie en 1597\footcite{LavabreBertrand2021Esprit}. La démonstration des simples consiste à identifier et présenter les substances médicinales, principalement végétales, à en décrire les propriétés et à apprendre leur emploi ; elle associe donc botanique, matière médicale et observation directe, notamment au Jardin des plantes\footnote{Fondé en 1593 par Henri IV et confié à Pierre Richer de Belleval, le Jardin des plantes de Montpellier est lié dès son origine à l'enseignement de l'anatomie et de la démonstration des simples\footcite{JardinPlantesMontpellier1984}. Il précède le Jardin royal des plantes médicinales de Paris, fondé en 1626 et ouvert au public en 1640\footcite{MNHNJardinPlantesParis}.}. Ces créations prolongent le contrôle revendiqué sur chirurgiens et apothicaires au nom de la santé publique\footcite[p. 321--322]{Irissou1934Pharmacie}. Visites de boutiques, examens professionnels et règlements étendent ainsi à l'époque moderne la certification universitaire de la compétence.

L'édit de 1610 limite enfin le nombre des docteurs agrégés participant au gouvernement et au partage des émoluments et organise leur remplacement\footcite[fol.~160v--164v]{AHFMS1}\footcite[p. 126--141]{Astruc1767}. L'intervention monarchique ne supprime pas les compétences académiques de la corporation mais garantit ses privilèges, finance les régences, spécialise les enseignements et encadre l'accès aux positions donnant part au pouvoir et aux revenus.

Cette évolution éclaire la composition des Archives historiques. Les actes pontificaux, seigneuriaux et royaux fondent, confirment ou réforment les compétences ; les statuts les traduisent en règles internes ; délibérations, comptes, nominations et actes de grades en enregistrent l'application. Le fonds conserve ainsi les niveaux documentaires successifs d'un même gouvernement, depuis l'octroi du droit jusqu'à son exercice par le collège et ses officiers.

\subsection{De la corporation médicale à la composante universitaire contemporaine}

À partir du XVIII\textsuperscript{e} siècle, l'histoire institutionnelle ne procède plus principalement par accumulation de privilèges. La Faculté diversifie ses enseignements, disparaît juridiquement pendant la Révolution, puis s'insère dans les organisations nationales établies par l'État. La permanence locale de la médecine recouvre donc plusieurs ruptures de statut, de tutelle et de production documentaire.

\subsubsection{Diversification des savoirs et des lieux d'enseignement au XVIII\textsuperscript{e} siècle}

Montpellier demeure l'un des principaux centres médicaux du royaume par sa réputation, ses effectifs et la densité de ses institutions sanitaires et savantes\footcite{Lacour2026GensSante}\footcite{Dulieu1958MouvementScientifique}. L'Université de médecine reste dominée par le collège des docteurs, mais coexiste avec chirurgiens, apothicaires, hôpitaux, Collège de chirurgie et sociétés savantes : la production et la transmission du savoir médical ne relèvent plus d'elle seule.

La médecine et la chirurgie étaient depuis le Moyen Âge des activités institutionnellement distinctes plutôt que deux branches d'une profession unifiée. Les médecins, formés dans les universités, revendiquaient une science fondée sur les textes et la théorie ; les chirurgiens, longtemps organisés en communautés professionnelles et formés par apprentissage, exerçaient un art manuel, même si certains recevaient une formation anatomique et savante. Cette séparation, variable selon les lieux et les périodes, se maintient sous l'Ancien Régime tout en devenant plus poreuse avec le développement des collèges de chirurgie, de l'anatomie et de la clinique\footcite{Brockliss1989EnseignementMedical}\footcite{Brier2019OrganisationMedicale}.

À Montpellier, la spécialisation amorcée à la fin du XVI\textsuperscript{e} siècle se poursuit par de nouveaux enseignements. La médecine pratique et le service des pauvres renforcent la place de l'observation des malades, même si la clinique demeure irrégulière et ne remplace pas le commentaire des textes. Cette évolution participe d'un mouvement européen de rapprochement entre théorie, observation et pratique hospitalière, antérieur à sa systématisation révolutionnaire\footcite{Brockliss1989EnseignementMedical}\footcite{Dulieu1955VieMedicale}.

L'ouverture de l'Hôtel Saint-Côme en 1757 donne au Collège de chirurgie des espaces d'enseignement et de démonstration anatomique. Les deux corps restent juridiquement séparés, mais leurs pratiques pédagogiques convergent. La Faculté affirme parallèlement une orientation doctrinale illustrée par le vitalisme de Théophile de Bordeu et de Paul-Joseph Barthez\footcite{LavabreBertrand2021Esprit}\footcite{Dulieu1971Barthez}.

Chaires, offices, examens et droits restent organisés dans le cadre corporatif, mais celui-ci s'est transformé avant la Révolution par la spécialisation des enseignements, le développement de la clinique, le rapprochement de la médecine et de la chirurgie et l'insertion dans un espace savant plus large. À la tête de l'ancienne Université demeure le chancelier, tandis que le doyen appartient au collège des docteurs sans exercer encore les fonctions de direction qui seront attachées à ce titre au XIX\textsuperscript{e} siècle. La rupture révolutionnaire met fin à cette organisation : avec la suppression de l'Université disparaît également la chancellerie qui en constituait l'un des principaux offices.

\subsubsection{De la suppression révolutionnaire à l'École de santé}

La législation révolutionnaire désorganise progressivement le cadre institutionnel de l'enseignement et de l'exercice médicaux. La loi des 2--17 mars 1791 supprime les maîtrises, jurandes, offices et privilèges professionnels, tout en proclamant la liberté d'exercer la profession de son choix ; le décret du 15 septembre 1793 prononce ensuite la suppression des facultés de médecine sur l'ensemble du territoire\footcite[art.~2 et 7]{LoiAllarde1791}\footcite[art.~3]{DecretFacultes1793}. À Montpellier, l'Université de médecine perd ainsi sa personnalité institutionnelle, ses offices et ses privilèges, sans que toute activité médicale cesse immédiatement : médecins et chirurgiens continuent d'enseigner et de répondre aux besoins de formation et de soins dans un cadre provisoire\footcite{Dulieu1955VieMedicale}\footcite{HuardImbaultHuart1973}.

La pénurie de praticiens, la désorganisation des formations et les besoins des armées conduisent la Convention à rétablir rapidement un enseignement public. Le décret du 14 frimaire an III, soit le 4 décembre 1794, crée trois Écoles de santé à Paris, Montpellier et Strasbourg\footcite[art.~I, p.~402]{DecretEcoles1794}\footcite{Brockliss1989EnseignementMedical}. Le choix de Montpellier reconnaît son importance médicale, mais la nouvelle école n'est pas la continuation juridique de l'ancienne Université : placée sous l'autorité de l'État, elle réunit médecine et chirurgie et reçoit un personnel, un programme et une organisation nouveaux. La direction de l'établissement relève désormais d'un directeur, fonction notamment exercée par Charles-Louis Dumas, nommé professeur de la nouvelle École de santé en 1794 et appelé à en assurer la direction\footcite{Condette2006Dumas}.

Avant 1795, l'Université ne demeure pas exclusivement dans les maisons de ses professeurs. Après une longue période durant laquelle les maîtres enseignent dans des locaux dispersés, notamment privés, elle utilise le Collège des Douze-Médecins et possède, à la fin de l'Ancien Régime, des bâtiments propres près de la rue du Bout-du-Monde, future rue de l'École-de-Pharmacie. Ces locaux, comprenant deux salles et un amphithéâtre d'anatomie installé dans une tour, sont devenus trop étroits et vétustes\footcite{Dulieu1955VieMedicale}\footcite{LavabreBertrand2021Esprit}.

En 1795, l'École de santé s'installe dans l'ancien monastère Saint-Benoît, devenu palais épiscopal puis bien national, aujourd'hui site historique de la Faculté. La continuité locale tient aux enseignants, aux pratiques et à certains documents ou objets réemployés ; la nouvelle institution doit néanmoins constituer ses propres instruments d'étude. Les collectes menées à Montpellier, Paris et dans les dépôts révolutionnaires, puis les missions de Gabriel Prunelle, forment une bibliothèque adaptée au nouveau programme, aux côtés du Jardin des plantes, du cabinet de physique et du Conservatoire d'anatomie\footcite{Denton2022Bibliotheque}. « Créer les ressources documentaires d'une institution nouvelle » signifie donc réunir, classer et rendre accessibles les livres, manuscrits, préparations anatomiques, instruments et spécimens nécessaires à l'enseignement réorganisé, et non seulement transférer intactes les collections de l'ancienne corporation.

La réforme articule davantage formation, dissection, observation clinique, hôpitaux et besoins sanitaires. En réunissant l'enseignement médical et chirurgical, elle contribue à effacer leur ancienne hiérarchie institutionnelle et à former une profession plus unifiée\footcite{HuardImbaultHuart1973}\footcite{Brockliss1989EnseignementMedical}. L'École de santé n'est donc pas l'Université restaurée, mais un établissement public nouveau qui reprend localement une partie de ses personnels, ressources et fonctions.

\subsubsection{La Faculté dans l'Université nationale}

La loi du 19 ventôse an XI, soit le 10 mars 1803, rétablit une réglementation nationale de l'exercice médical et distingue docteurs en médecine ou en chirurgie et officiers de santé\footcite{LoiVentoseAnXI}. Les Écoles de santé deviennent Écoles de médecine, puis la réorganisation de l'enseignement supérieur engagée sous l'Empire les intègre au nouvel ordre universitaire. À Montpellier, Charles-Louis Dumas, qui avait dirigé l'École, prend le titre de doyen de la Faculté de médecine ; le décret du 17 mars 1808 inscrit celle-ci dans l'Université impériale\footcite{DecretUniversiteImperiale1808}\footcite{Condette2006Dumas}. Le changement de titre accompagne ainsi une transformation beaucoup plus profonde de l'institution : le doyen de la Faculté du XIX\textsuperscript{e} siècle ne succède pas au chancelier comme titulaire du même office, mais prend place dans une organisation nouvelle, centralisée et hiérarchisée par l'État.

La Faculté ne restaure donc pas la corporation d'Ancien Régime. Elle constitue désormais un établissement spécialisé d'un système national dans lequel programmes, grades et nominations relèvent de l'État. Le doyen en assure la direction à l'échelle de la Faculté, sous l'autorité des instances de l'Université et de l'administration centrale ; le titre acquiert ainsi une fonction institutionnelle différente de celle qu'il possédait dans l'ancienne Université de médecine. Les diplômes ne donnent plus accès à une communauté autonome de maîtres, mais sanctionnent un cursus national et conditionnent l'exercice professionnel. Ce modèle se maintient au XIX\textsuperscript{e} siècle malgré les différences entre docteurs, officiers de santé et autres catégories de praticiens\footcite{Leonard1966EtudesMedicales}\footcite{Brier2019OrganisationMedicale}.

Les facultés de médecine, sciences, lettres et droit restent longtemps des établissements distincts. La création d'un conseil général des facultés en 1885 amorce leur rapprochement ; la loi du 10 juillet 1896 appelle \enquote{université} le corps formé, dans chaque académie, par la réunion des facultés de l'État\footcite{LoiUniversites1896}. L'Université de Montpellier est ainsi reconstituée comme fédération, au sein de laquelle la médecine conserve ses bâtiments, ses enseignements, ses collections et son administration.

\subsubsection{Les recompositions universitaires depuis 1968}

La loi d'orientation du 12 novembre 1968, dite loi Faure, remplace les facultés par des unités d'enseignement et de recherche et réorganise les universités selon les principes d'autonomie, de participation et de pluridisciplinarité\footcite{LoiFaure1968}. À Montpellier, son application conduit en 1970 à trois établissements : Montpellier 1 réunit principalement droit, économie, gestion, médecine, odontologie et pharmacie ; Montpellier 2 rassemble sciences et techniques ; Montpellier 3, devenue Université Paul-Valéry, regroupe lettres, langues, arts et sciences humaines et sociales\footcite{CNE2004Montpellier}. En 1984, la loi du 26 janvier sur l'enseignement supérieur, dite loi Savary, remplace les unités d'enseignement et de recherche par des unités de formation et de recherche, sans faire disparaître dans l'usage le terme de \enquote{faculté}\footcite[art.~25 et 32]{LoiSavary1984}.

Cette organisation répartit la production documentaire entre plusieurs niveaux. L'UFR produit et conserve notamment les dossiers liés aux enseignements, examens, stages, instances et activités propres à la médecine ; la présidence et les services centraux traitent les décisions communes, budgets, marchés, immobilier, carrières et politiques d'établissement ; les services partagés, tels que bibliothèques, archives, ressources humaines ou systèmes d'information, produisent leurs propres dossiers. Un même projet peut donc laisser des traces complémentaires dans plusieurs services, ce qui explique l'éclatement contemporain des responsabilités de conservation.

Le 1\textsuperscript{er} janvier 2015, Montpellier 1 et Montpellier 2 fusionnent pour former l'actuelle Université de Montpellier, tandis que Paul-Valéry demeure distincte\footcite{DecretUniversiteMontpellier2014}. La Faculté de médecine devient une composante du nouvel établissement, tout en conservant son identité, ses implantations et ses services. Inauguré en 2017, le campus Arnaud-de-Villeneuve accueille au nord de la ville les enseignements médicaux dans des locaux adaptés aux effectifs contemporains et aux équipements de pointe ; il complète le bâtiment historique, qui conserve des fonctions institutionnelles, patrimoniales et pédagogiques\footcite{FaculteMedecineCampusADV2017}.

La continuité revendiquée de l'enseignement médical ne doit donc pas être confondue avec l'identité juridique de l'institution. Écoles de maîtres, \emph{universitas medicorum}, Faculté d'Ancien Régime, École de santé, École puis Faculté de médecine impériale, faculté, UER puis UFR composante de Montpellier 1 et enfin de l'Université de Montpellier constituent des organisations successives. Ces transformations déterminent les formes de production documentaire ainsi que les continuités et les ruptures qui structurent aujourd'hui les Archives historiques.

\chapter{Histoire et constitution des archives historiques de la Faculté}
\chapitrecourt{Histoire et constitution des archives}

\section{La mise en commun progressive des archives universitaires (1220--1583)}
\label{sec:mise-en-commun-archives}

Les documents présentés au chapitre précédent n'ont pas formé d'emblée un fonds centralisé. Statuts, privilèges, registres, actes de grades, comptes et pièces judiciaires demeurent d'abord auprès des autorités ou des officiers qui les produisent, les reçoivent et les utilisent. Leur conservation répond moins à une finalité patrimoniale qu'aux nécessités du gouvernement : appliquer une norme, vérifier un cursus, administrer des biens ou défendre un droit. L'histoire des archives commence ainsi par une garde distribuée, progressivement complétée par des dispositifs collectifs de sécurité.

\subsection{Une garde attachée aux fonctions}

Les statuts de 1220, premier texte normatif connu de l'Université de médecine, illustrent cette logique. Leur lecture lors de la réception des maîtres exige qu'ils demeurent disponibles ; trois exemplaires doivent donc être conservés par l'évêque de Maguelone, le prieur de Saint-Firmin et le chancelier\footcite[fol.~187--188, vues~186--187/348]{ADHG1126}\footcite{Conrad1220}\footcite[p.~IX, vue~28 de la version numérisée]{CalmetteCartulaire1912}. Cette pluralité ne constitue pas un fonds, mais garantit la transmission et l'opposabilité de la norme en limitant les conséquences de la perte d'un exemplaire.

Le plus ancien document actuellement conservé par les Archives historiques de la Faculté est l'original des statuts complémentaires de 1240\footcite{AHFMB1}\footcite{Statuts1240}. La distinction entre le premier acte institutionnel connu et la plus ancienne pièce du fonds actuel montre que l'histoire de l'Université et celle de ses archives ne se confondent pas. Les procédures anciennes d'inscription et de collation des grades supposent également des documents d'enregistrement aujourd'hui disparus : leur absence traduit vraisemblablement une rupture de transmission plutôt qu'une absence de production\footcite[p.~IX, vue~28 de la version numérisée]{CalmetteCartulaire1912}.

Aucun office archivistique autonome n'est alors identifié. La garde des actes découle des compétences exercées. En 1240, le chancelier jure ainsi de transmettre fidèlement les statuts à ses successeurs (\emph{statuta Universitatis fideliter suis successoribus reservaturum})\footcite{AHFMB1}\footcite{Statuts1240}. Vers 1390, étudiants et bacheliers précisent que les originaux des actes pontificaux et des légats sont conservés auprès de lui (\emph{quorum originalia apud cancellarium dicte Universitatis reservantur})\footcite{RequeteEtudiants1390}. La chancellerie détient donc les principaux titres nécessaires à la défense des droits universitaires, sans réunir l'ensemble des archives.

La décision du cardinal Jean de Saint-Marc de 1364 révèle au contraire leur dispersion. Chargé d'examiner les statuts et les usages de la Faculté, le commissaire pontifical ordonne à tous les détenteurs de droits, lettres, actes, privilèges, chartes ou écritures, publiques comme privées, de les lui remettre sur réquisition\footcite[fol.~8~v°, vues~15--16/194]{ADHG1277}\footcite{JeanSaintMarc1364}. Leur rassemblement ponctuel devient nécessaire pour instruire un différend, précisément parce que leur garde demeure répartie entre les membres de la communauté.

Cette organisation brouille la frontière entre archives institutionnelles et papiers personnels. Les documents appartiennent à l'Université, mais suivent les officiers qui les utilisent et peuvent demeurer dans leurs maisons ou leurs successions. Elle en facilite l'usage immédiat, au prix de risques de rétention, de confusion et de perte.

\subsection{Sécuriser les titres sans empêcher leur circulation}

Des dispositifs collectifs réduisent progressivement ces risques. Les statuts de 1340 placent les biens communs dans une \emph{archa}, conservée chez le plus ancien procureur et fermée par trois serrures dont les clés sont confiées à des maîtres différents\footcite{Statuts1340}. Le texte concerne d'abord les ressources financières et ne permet pas d'affirmer que ce coffre contient déjà les archives. Il établit néanmoins un principe de garde collégiale empêchant qu'un seul officier puisse disposer des biens communs.

En février 1474, ce principe est explicitement appliqué à la matricule. Le registre doit être déposé dans l'église Saint-Mathieu, dans une caisse fixée au mur par des chaînes ou des crampons de fer et fermée par trois clés distinctes. Celles-ci sont confiées respectivement au chancelier, au plus jeune maître et, alternativement d'une année à l'autre, à un bachelier ou à un écolier exerçant la fonction de procureur\footcite[p.~X, vue~29 de la version numérisée]{CalmetteCartulaire1912}. Cette répartition soumet l'ouverture de la caisse à l'accord de plusieurs membres de l'Université et renouvelle périodiquement une partie de ses détenteurs. Elle correspond à la valeur probatoire de la matricule, qui établit l'appartenance des étudiants à l'Université et permet de vérifier leur cursus.

Cette sécurisation n'interdit pas la circulation des titres. Les actes doivent être consultés, copiés et présentés aux autorités chargées de les appliquer. En juin 1396, le chancelier Jean Piscis remet au bailli de Montpellier les lettres de Charles VI autorisant la délivrance annuelle d'un corps de supplicié pour l'enseignement anatomique, puis fait dresser un instrument public constatant leur présentation et la réponse reçue\footcite{AHFMB17}\footcite{PresentationCadavres1396}. Le déplacement du titre produit ainsi un nouvel acte qui conserve la preuve de son usage.

Cette circulation concerne aussi les communautés professionnelles placées sous la juridiction universitaire. Les compagnons chirurgiens demandent au chancelier Laurent Joubert d'ordonner la restitution de leurs statuts et privilèges, retenus depuis deux ans par les maîtres chirurgiens, afin d'organiser leur étude et d'élire leur abbé\footnote{Le titre d'\enquote{abbé} ne désigne pas ici une dignité religieuse, mais un office électif au sein de la communauté des compagnons, notamment chargé de leur représentation.}\footcite{AHFML2}. La détention matérielle de ces textes conditionne donc la capacité du groupe à connaître ses règles, à désigner son représentant et à exercer ses droits.

Copies authentiques, instruments publics, récépissés et demandes de restitution répondent ainsi à une tension entre disponibilité et sécurité. Les archives ne sont pas encore conçues comme un ensemble immobile : leur valeur procède de leur mobilisation dans l'action administrative et judiciaire, tandis que leurs déplacements rendent nécessaire un contrôle croissant de leur garde.

\subsection{Le retour des documents à la garde commune}

Au XVI\textsuperscript{e} siècle, la dispersion devient un problème explicitement traité par l'assemblée. Le 1\textsuperscript{er} août 1571, les docteurs décident de récupérer \enquote{tous les registres publics et tous les actes} du Collège royal des médecins (\emph{omnes tabulas publicas omniaque instrumenta}) auprès du chancelier Antoine Saporta, du doyen Blezin, des héritiers de Guillaume Rondelet et de toute autre personne susceptible de les détenir (voir annexe, fig.~\ref{fig:s8-f51v})\footcite[fol.~51~v°]{AHFMS8}\footcite[p.~X, où Joseph Calmette donne la transcription du passage, vue~29 de la version numérisée]{CalmetteCartulaire1912}. Ils prescrivent leur dépôt dans le coffre public de la salle du Collège (\emph{in arcam publicam conclavis praedicti Collegii}). Dans cette délibération, le terme \emph{archa} désigne donc explicitement le contenant destiné aux archives de l'Université, et non plus seulement un coffre réservé à ses biens ou à ses ressources.

Le \emph{Collegium} désigne ici le Collège royal des médecins, c'est-à-dire le corps des docteurs régents constituant la Faculté, et non le Collège des Douze-Médecins fondé par Urbain V pour loger des étudiants boursiers. La salle du Collège désigne donc le local commun dans lequel les docteurs se réunissent et conservent certains biens de la corporation ; le texte ne permet toutefois pas de situer plus précisément le coffre en 1571.

La réunion des documents dans un coffre commun ne supprime pas la nécessité d'en confier la responsabilité à un membre de l'Université. Deux pièces conservées dans le fonds ancien indiquent que les archives doivent en principe être placées sous la garde du doyen ou du plus ancien docteur (voir annexe, fig.~\ref{fig:c9}, \ref{fig:c35-01} et~\ref{fig:c35-02})\footcite{AHFMC9}\footcite{AHFMC35}. Cette responsabilité demeure attachée à une fonction ou à l'ancienneté dans le collège, sans constituer encore un office archivistique spécialisé.

La propriété collective des documents est ainsi opposée à leur détention individuelle. En 1572, Laurent Joubert doit à son tour restituer les papiers de l'Université qu'il conserve (voir annexe, fig.~\ref{fig:c17})\footcite{AHFMC17}\footcite[p.~XI--XII]{CalmetteCartulaire1912}. D'autres documents, retrouvés dans la succession de Pierre Trémolet après sa mort à Paris, avaient déjà été rendus à la Faculté (voir annexe, fig.~\ref{fig:c4-recto})\footcite{AHFMC4}\footcite[p.~XII, vue~31 de la version numérisée]{CalmetteCartulaire1912}.

Le coffre de 1571 ne constitue pas encore un dépôt au sens contemporain, faute de service spécialisé, de versements réguliers et d'un classement systématique. Les documents continuent d'en sortir selon les besoins. Le 11 février 1576, le procureur Guichard reconnaît ainsi avoir reçu du chancelier l'une des trois clés du coffre, en qualité de plus jeune docteur, puis du doyen Hucher le livre des privilèges afin de le présenter (voir annexe, fig.~\ref{fig:s8-f90v})\footcite[fol.~90~v°]{AHFMS8}. À partir de cette mention, Calmette suppose que Guichard remplit alors les fonctions d'archiviste de l'Université\footcite[p.~XI, vue ~30 de la version numérisée]{CalmetteCartulaire1912}. Il ne s'agit pas nécessairement d'un office constitué sous ce titre, mais d'une responsabilité pratique portant sur la garde, la remise et le suivi des documents. Le récépissé concilie ainsi la circulation de la pièce avec l'identification de son détenteur.

Les recherches entreprises depuis 1571 ont vraisemblablement sauvé des documents qui auraient pu disparaître dans les papiers de leurs détenteurs\footcite[p.~XI--XIII, vues~30 à 32 de la version numérisée]{CalmetteCartulaire1912}. Elles préparent surtout une transformation décisive : après avoir réaffirmé la propriété collective des archives et tenté de les réunir, la Faculté entreprend de les identifier et de les classer.

\section{L'inventaire de 1583 : la constitution d'un ordre documentaire}

Le premier inventaire général connu des archives de la Faculté est établi en 1583 sous l'autorité du chancelier Jean Hucher, assisté du doyen Nicolas Dortoman et du docteur régent Jean Saporta. Son préambule indique que les archives susceptibles d'être \enquote{recouvert des egarez} ont été recherchées avant leur mise en ordre\footcite[p.~XII, vue~31 de la version numérisée]{CalmetteCartulaire1912}. L'opération prolonge donc les restitutions des années 1570.

L'analyse qui suit repose principalement sur l'édition donnée par Joseph Calmette dans le second tome du \emph{Cartulaire de l'Université de Montpellier}, qui reproduit l'inventaire et en présente la structure\footcite[p.~XXII--XXIX et p.~1--206, vues~41 à 51 et 187 à 394 de la version numérisée]{CalmetteCartulaire1912}.

\subsection{Un classement fondé sur les fonctions de la Faculté}

Les documents sont répartis en quatorze sacs correspondant aux principales affaires de l'Université : privilèges, gages et pensions, anatomie, graduations, régences, arrêts, procès, comptes, formulaires, livres et relations avec les chirurgiens\footcite[p.~XXV, vue~47 de la version numérisé]{CalmetteCartulaire1912}. Ce classement n'est ni chronologique ni fondé sur le support des documents. Il regroupe les actes selon les fonctions auxquelles ils se rapportent et offre ainsi une représentation documentaire de l'institution.

Les privilèges sont isolés afin de pouvoir être produits rapidement en cas de contestation ; les dossiers de graduations, de régences et d'anatomie documentent les activités académiques ; comptes, pensions et procédures répondent aux besoins administratifs et contentieux. L'inventaire doit rendre les titres disponibles pour gouverner et défendre l'Université.

Chaque sac est subdivisé en liasses désignées par des lettres, redoublées ou triplées lorsque l'alphabet ne suffit plus, les pièces sont ensuite distinguées par une lettre minuscule ou un numéro\footcite[p.~XXV--XXVI, vues~47 à 48 de la version numérisée]{CalmetteCartulaire1912}. Les cotes, reportées sur les documents, relient leur emplacement matériel à leur description : une mention comme \enquote{sac I C} renvoie simultanément au contenant, à la liasse et à l'inventaire\footcite[p.~XXVI, vue~48 de la version numérisée]{CalmetteCartulaire1912}.

Les notices indiquent généralement la nature, l'auteur, la date et l'objet de l'acte\footcite[p.~XXVIII--XXIX, vues~50 à 51 de la version numérisée]{CalmetteCartulaire1912}. Cotation et analyse permettent donc de localiser une pièce et d'en apprécier le contenu avant de la consulter. L'inventaire constitue à la fois un état du fonds, un instrument de recherche et un moyen de contrôler la présence des documents.

\subsection{Un instrument tenu à jour}

L'inventaire reste utilisé après 1583. Ses additions marginales signalent de nouvelles pièces, complètent des analyses et corrigent certaines descriptions\footcite[p.~XIII, XXII--XXIII, vues~32 et 41 à 43 de la version numérisée]{CalmetteCartulaire1912}. La structure en sacs, liasses et articles permet ces accroissements sans refonte générale.

Les archives pourraient alors avoir été placées sous la garde du secrétaire de la Faculté\footcite[p.~XIII, vue~32 de la version numérisée]{CalmetteCartulaire1912}. Cette fonction est alors exercée par un notaire choisi par la Faculté. Théodore Dégan occupe ainsi le secrétariat sous les chanceliers Laurent Joubert et Jean Hucher ; après sa mort, l'assemblée du 18 février 1593 désigne pour lui succéder \enquote{maistre Pierre Gallet, notaire royal de la ville de Montpellier}\footcite[fol.~184~v\textsuperscript{o}]{AHFMS8}. Le rapprochement entre le secrétariat et la garde des archives est également attesté dans d'autres institutions universitaires de la Renaissance, où la production et la tenue des registres sont progressivement confiées à un secrétaire ou à un officier spécialisé\footcite{PeekHall1962Cambridge}. Sa responsabilité porte d'abord sur la rédaction et l'authentification des actes nécessaires au fonctionnement de l'institution, auxquelles peut s'ajouter la conservation des documents ainsi produits.

L'entreprise de 1583 marque le passage d'une conservation principalement attachée aux détenteurs à une maîtrise institutionnelle fondée sur trois opérations indissociables : rassemblement, classement et description. Elle n'immobilise pas les archives, toujours utilisées dans l'administration et les procès, mais fournit un cadre permettant d'en contrôler les mouvements.

Cette évolution s'inscrit dans une transformation plus générale des pratiques européennes de gouvernement entre la fin du Moyen Âge et l'époque moderne. L'accroissement de la production écrite conduit alors de nombreuses institutions à différencier les espaces de conservation, à multiplier inventaires et systèmes de cotation et à organiser les documents comme des ressources probatoires et informationnelles\footcite{Head2016Archives}\footcite{Head2019MakingArchives}. L'inventaire montpelliérain relève de ce mouvement sans constituer encore un système archivistique moderne : son efficacité demeure tributaire de la continuité de sa mise à jour et de la régularité des réintégrations.

\section{Usages, pertes et reclassement aux XVII\textsuperscript{e} et XVIII\textsuperscript{e} siècles}

\subsection{L'érosion du classement de 1583}

Les additions à l'inventaire se raréfient au cours du XVII\textsuperscript{e} siècle, tandis que les nouvelles pièces ne sont plus systématiquement enregistrées\footcite[p.~XIII, vue~32 de la version numérisée]{CalmetteCartulaire1912}. L'instrument cesse progressivement de correspondre à l'état réel du fonds.

Cette rupture est aggravée par l'usage intensif des archives. Registres et titres sont empruntés par les professeurs, mobilisés dans les litiges internes ou dans les conflits avec chirurgiens et apothicaires, puis transportés devant les juridictions. Certaines pièces quittent notamment Montpellier pour être communiquées aux avocats et procureurs chargés des affaires de la Faculté devant le Parlement de Toulouse. Trois lettres attestent ainsi la réception de registres d'inscription à Toulouse en 1708, 1715 et 1716 (voir annexe, fig.~\ref{fig:f24-01})\footcite{AHFMF22}\footcite{AHFMF23}\footcite{AHFMF24}. Ces accusés de réception permettent de suivre ponctuellement la circulation des documents, sans toutefois garantir leur retour effectif dans les archives de la Faculté\footcite[p.~XIV, vue~33 de la version numérisée]{CalmetteCartulaire1912}.

Les procédures jugées à Toulouse peuvent également engendrer de nouvelles copies ou transcriptions conservées à Montpellier. Le registre des \emph{Arrêts et Déclarations} réunit ainsi plusieurs décisions du Parlement, dont l'arrêt du 14 mars 1651 annulant deux agrégations et ordonnant leur mise au concours\footcite[fol.~85~v\textsuperscript{o}, acte~91]{AHFMS3}. Il ne constitue pas une preuve directe du transport de l'original, mais montre comment l'activité contentieuse alimente en retour les archives de la Faculté.

Les documents déplacés ne retrouvent pas toujours leur emplacement initial et certains disparaissent. Ces difficultés sont accentuées par l'incommodité du rangement en sacs et par la complexité des cotes portées au dos des pièces, qui prête à confusion\footcite[p.~XIII, vue~42 de la version numérisée]{CalmetteCartulaire1912}. Le classement de 1583 reste efficace tant que ces cotes sont respectées, mais devient difficile à rétablir dès que les réintégrations se relâchent\footcite[p.~XIII--XIV, vues~42 à 43 de la version numérisée]{CalmetteCartulaire1912}. Le désordre procède de l'écart croissant entre un inventaire qui n'est plus actualisé et un fonds continuellement accru, consulté et déplacé.

Lorsqu'il examine les archives au XVIII\textsuperscript{e} siècle, Jean-François Imbert constate qu'une lacune de l'inventaire de 1583 ne résulte pas d'une simple négligence. Le titre du sac~VI, principalement consacré aux pièces relatives à la chancellerie, a été gratté dans le manuscrit. Sur une fiche ajoutée à l'inventaire, Imbert indique que ce sac renfermait \enquote{tout ce qui avoit rapport à la charge de chancelier}, qu'il \enquote{a été enlevé des archives il y a longtemps} et que les feuillets de l'inventaire général \enquote{où étoient nottées et cottées les matières y relatives ont été arrachés}.\footnote{La fiche porte exactement la mention suivante : \enquote{Ce sac, qui renfermoit tout ce qui avoit rapport a la charge de Chancelier, a été enlevé des archives il y a longtemps ; et deux feuillets de l'inventaire general ou etoient nottées et cottées les matieres y relatives ont été arrachés. On y a suppléé en ramassant, autant qu'il a été possible, des copies et extraictz concernant les droits du Chancelier. Le tout est au coffre 4\textsuperscript{e} de mes recueils ou l'on trouvera les pieces} \parencite[p.~XXIV--XXV, vues~46 à 47 de la version numérisée]{CalmetteCartulaire1912}.} Il tente de suppléer cette perte en réunissant des copies, des extraits et divers documents relatifs aux droits du chancelier\footcite{AHFMQ1}. La disparition conjointe du dossier et des notices qui permettaient de l'identifier témoigne d'une soustraction volontaire, sans qu'il soit possible d'en déterminer la date ni l'auteur. Ce constat éclaire l'état du fonds auquel Imbert est confronté lorsqu'il entreprend sa réorganisation en 1769.

\subsection{La réorganisation conduite par Jean-François Imbert au XVIII\textsuperscript{e} siècle}

En 1769, le procureur général près le Parlement de Toulouse rappelle à la Faculté son obligation de maintenir les archives dans un état qui permette leur consultation. Il relève que plusieurs registres sont conservés depuis des années par des particuliers, au risque de compromettre les droits de l'Université (voir annexe, fig.~\ref{fig:f54-19} et~\ref{fig:f54-20}).\footnote{\enquote{Il est revenu, Monsieur, à M. le Procureur Général, qu'il se commet journellement bien des abus dans la manière dont sont tenus les registres de la Faculté de Médecine de la ville de Montpellier ; que ces registres, qui doivent être déposés dans les Archives de cette Faculté ou dans les mains du secrétaire, sont toujours entre les mains des professeurs, et que ces infractions aux règles ont occasionné l'égarement et la perte de plusieurs registres} \parencite{AHFMF54}.} La dispersion des papiers affecte la capacité de l'institution à administrer ses affaires et à produire ses titres.

Le chancelier Jean-François Imbert entreprend alors la récupération et le reclassement des documents ainsi que la rédaction d'un nouvel inventaire\footcite[p.~XV--XVII, vues~34 à 35 de la version numérisée]{CalmetteCartulaire1912}\footcite{Dulieu1956Imbert}. Les pièces préparatoires conservées témoignent de l'ampleur de cette opération\footcite{AHFMA5}. La décision de procéder à ce nouvel inventaire est consignée dans le registre des délibérations de la Faculté (voir annexe, fig.~\ref{fig:s17-p8}).\footnote{Calmette date cette décision du 12 septembre 1769, mais la délibération correspondante est en réalité inscrite à la date du 1\textsuperscript{er} juillet 1769 dans le registre de la Faculté \parencite{AHFMS17}. Le 12 septembre correspond au commencement de la transcription des arrêts, déclarations et autres titres dans un registre distinct \parencite[fol.~1]{AHFMS3}.}

L'ordre matériel n'est plus fondé sur les quatorze sacs de 1583. Les pièces sont désormais réparties dans des armoires, elles-mêmes subdivisées en compartiments, liasses ou dossiers auxquels renvoient de nouvelles cotes. Le contenant fixe remplace ainsi le sac mobile, tandis que l'inventaire décrit la localisation des documents dans le mobilier. Cette réorganisation répond à l'accroissement du fonds, mais superpose un nouveau système aux anciennes cotes encore inscrites sur les pièces. Elle ne restitue donc pas l'ordre de 1583 : elle en conserve certains regroupements tout en les adaptant à un rangement devenu plus stable.

Imbert fait également transcrire dans un registre les arrêts, déclarations et autres titres jugés importants pour la Faculté. Commencé le 12 septembre 1769 et poursuivi jusque vers 1780, ce recueil compte 239 actes classés en catégories approximatives plutôt qu'en stricte succession chronologique (voir annexe, fig.~\ref{fig:s3-f1}).\footnote{Le premier feuillet porte la déclaration suivante : \enquote{Nous Chancelier et professieur[s] soussignés, chargés par l'Université de faire et dresser l'inventaire des titres et pièces renfermés dans les archives, avons jugé à propos de faire transcrire au présent registre les arrêts, déclarations, et autres titres qu'on trouvera cy après, concernant l'Université à Montpellier. Le 12 septembre 1769. Imbert, chancelier et juge. René ; Broussonet} \parencite[fol.~1]{AHFMS3}. Commencée en 1769, la transcription est achevée vers 1780. Jusqu'au folio~38, les pièces transcrites sont suivies de l'apostille : \enquote{Collationné [...] par moy, Secrétaire de l'Université de Médecine de Montpellier. Vincent, secrétaire} \parencite[p.~XXIII, vue~48 de la version numérisée]{CalmetteCartulaire1912}.} La copie complète ainsi le classement matériel en réunissant dans un même volume le contenu des principaux titres.

L'opération ne répare pas les pertes antérieures. Plusieurs pièces inscrites dans l'inventaire de 1583 demeurent introuvables et certains registres ne peuvent être récupérés\footcite[p.~XVII--XVIII]{CalmetteCartulaire1912}. Elle constitue néanmoins une nouvelle étape de stabilisation : comme dans les années 1570, la récupération des documents précède leur réorganisation. La répétition de cette séquence montre que l'ordre archivistique doit être périodiquement rétabli à mesure que les usages et les accroissements distendent le lien entre les archives, leur classement et leur inventaire.

\section{Le reclassement de Joseph Calmette et la patrimonialisation du fonds ancien (1903--1912)}

La suppression de l'Université de médecine à la Révolution puis les réorganisations du XIX\textsuperscript{e} siècle modifient le statut et les usages de ses archives. Au début du XX\textsuperscript{e} siècle, celles-ci ne servent plus au gouvernement de la corporation disparue : elles constituent des sources historiques dont l'identification suppose un nouveau traitement.

\subsection{Un nouvel inventaire à partir de la reconstitution des fonds}

Le projet de classement est engagé par une délibération du conseil de la Faculté du 24 juillet 1902 (voir annexe, fig.~\ref{fig:1med51-calmette})\footcite{AHFM1MED51}. En 1903, la Faculté de médecine et le Conseil de l'Université confient à l'archiviste-paléographe Joseph Calmette le classement des archives anciennes. Lorsqu'il commence son travail, les registres sont dispersés, les pièces médiévales et modernes mêlées à des dossiers du XIX\textsuperscript{e} siècle et les archives de l'ancienne Faculté confondues avec celles du Collège royal de chirurgie. Environ vingt-cinq mille documents, réunis sans ordre stable en liasses ou en paquets, doivent être examinés individuellement. L'opération, prévue pour cinq ans, est achevée en 1908\footcite[p.~I, vue~20 de la version numérisée]{CalmetteCartulaire1912}.

Calmette dépouille chaque pièce, l'identifie et la rapproche provisoirement de documents de même provenance ou de même fonction avant de constituer les séries. Son cadre procède de l'analyse progressive de la composition du fonds. Il cherche à concilier l'identification des institutions productrices avec la commodité de consultation, tout en reconnaissant que le chevauchement des compétences et la diversité des affaires imposent des choix\footcite[p. I-II, vues~20 à 21 de la version numérisée]{CalmetteCartulaire1912}.

Le classement sépare d'abord les deux fonds jusque-là mêlés. Les archives de l'ancienne Faculté de médecine occupent les séries A à G, consacrées à l'administration et aux privilèges, aux statuts, aux assemblées, à la comptabilité, aux grades, aux biens et fondations, puis à la correspondance et aux pièces diverses. Celles du Collège royal de chirurgie sont réparties dans les séries H à O selon une organisation comparable. Des séries particulières accueillent registres, cartulaires, inventaires, manuscrits et collections qui ne peuvent être rattachés sans artifice à l'un de ces ensembles\footcite[p.~XXX-XXXII, vues~52 à 54 de la version numérisée, suivies de notices sur les différentes séries et de transcriptions]{CalmetteCartulaire1912}.

Le travail de Calmette intervient au moment où les principes de provenance et de respect des fonds s'imposent progressivement dans la pratique archivistique européenne. Le \emph{Manuel pour le classement et la description des archives} de Samuel Muller, Johan Feith et Robert Fruin, publié en néerlandais en 1898 puis adapté en français en 1910, définit le fonds comme un ensemble organique résultant de l'activité de son producteur et recommande de fonder son classement sur l'organisation de celui-ci\footcite{MullerFeithFruin1910}\footcite{Duchein1977RespectFonds}. Il n'est pas nécessaire de supposer une application directe de ce manuel pour constater que la séparation opérée par Calmette entre Faculté de médecine et Collège de chirurgie relève de la même culture professionnelle.

Son cadre alphabétique substitue ainsi à la localisation en sacs ou en armoires une organisation fondée sur les producteurs et leurs fonctions. Il ne restitue pas matériellement les classements antérieurs, mais en conserve la mémoire grâce aux anciennes cotes relevées sur les pièces.

\subsection{Le tome II du \emph{Cartulaire de l'Université de Montpellier}}

Dès le 9 novembre 1907, le conseil de la Faculté remercie Calmette pour son travail et décide la publication du \emph{Cartulaire} (voir annexe, fig.~\ref{fig:1med52-calmette})\footcite{AHFM1MED52}. Le traitement aboutit en 1912 à la publication du deuxième volume du \emph{Cartulaire de l'Université de Montpellier}. À la différence du premier volume de 1890, consacré à l'histoire de l'institution et à l'édition de ses principaux actes médiévaux, celui-ci contient l'inventaire numérique des archives anciennes jusqu'en 1808, accompagné de notices historiques et diplomatiques\footcite{Fournier1913Cartulaire}\footcite{Galabert1914Cartulaire}.

Calmette retrouve au cours du classement l'inventaire de 1583, tenu pour perdu lors de la publication du premier volume\footcite[p.~XXXII, vue~59 de la version numérisée]{GermainUniversite1890}. Il en donne une transcription complète, suppléée pour certains folios manquants par un inventaire sommaire contemporain\footcite{AHFMA2}. Il publie également plusieurs actes redécouverts\footcite[p.~851-869, vues~1039 à 1057 de la version numérisée]{CalmetteCartulaire1912}. Un tableau de concordance relie les cotes de 1583 au nouveau classement\footcite[p.~CXLVII-CLVIII, vues~173 à 184 de la version numérisée]{CalmetteCartulaire1912}, signale les identifications incertaines et rend visibles les pertes intervenues entre les campagnes successives.

L'instrument de 1912 permet ainsi d'étudier simultanément le fonds ancien et sa transmission. Les anciennes cotes, les pièces manquantes et les rapprochements proposés par Calmette documentent les états successifs des archives autant que leur contenu. Son cadre de classement demeure aujourd'hui la référence pour le fonds antérieur à 1808, dont l'Université distingue officiellement les archives modernes et contemporaines encore en cours de traitement\footcite{UMPatrimoineArchives}.

Le reclassement modifie enfin le statut des documents. Installés au Musée Atger dans un mobilier conçu pour leur conservation, certains sont exposés dans une vitrine centrale formant, selon Calmette, un \enquote{musée des archives}\footcite[p.~XVI-XXI, vues~35 à 40 de la version numérisée]{CalmetteCartulaire1912}. Désormais soustrait aux besoins administratifs courants, le fonds ancien devient un patrimoine universitaire à conserver, étudier et présenter. Cette patrimonialisation rejoint plus largement la reconnaissance progressive des collections et des archives des établissements d'enseignement supérieur\footcite{MadayMechine2011}\footcite{Hottin2008}.

\section{Le traitement des archives modernes et contemporaines aux XX\textsuperscript{e} et XXI\textsuperscript{e} siècles}

La borne de 1808 retenue par Calmette correspond au passage de l'ancienne corporation à la Faculté intégrée à l'Université impériale, présenté au chapitre précédent. Les documents postérieurs forment dès lors un ensemble moderne et contemporain dont le traitement suit une histoire distincte.\footnote{Cette section n'aurait pu être rédigée sans l'aide de Sophie Dikoff, qui a mené des recherches sur l'histoire des archives modernes et contemporaines, notamment dans le cadre de l'élaboration de la feuille de route 2025--2029 présentée plus loin dans ce mémoire.}

\subsection{La campagne de classement sous le décanat de Gaston Giraud}

Durant la Seconde Guerre mondiale, le fonds ancien est mis à l'abri à Saint-Guilhem-le-Désert\footcite{AHFM1MED211}. Yvonne Vidal, conservatrice de la bibliothèque, en poursuit ensuite l'étude et publie en 1958 une présentation de la bibliothèque et des archives de la Faculté\footcite{Vidal1958BibliothequeArchives}\footnote{Ses carnets de notes, conservés à la BUHM depuis leur redécouverte dans les archives de la bibliothèque lors du déménagement des collections de 2021, indiquent qu'elle a travaillé sur les archives durant la guerre.}. Son travail participe au maintien d'une connaissance commune des collections imprimées, manuscrites et archivistiques alors conservées dans les mêmes espaces.

Après le conflit, Gaston Giraud, doyen de 1941 à 1960, engage une vaste modernisation des locaux. Les archives conservées dans le fonds Giraud\footnote{Complémentaires du fonds de l'architecte Jean de Richemond conservé aux Archives départementales de l'Hérault.} documentent les travaux du bâtiment historique et de l'Institut de biologie, l'aménagement des salles établies dans les anciens mâchicoulis et, plus largement, l'extension des sites de la Faculté\footnote{Elles ont notamment servi, en 2021, à la préparation de l'exposition \emph{Les Bâtisseurs de la Faculté de médecine du XX\textsuperscript{e} siècle} \parencite{BatisseursFaculte2021}.}.

La réorganisation matérielle s'accompagne d'une reprise des archives du XIX\textsuperscript{e} siècle, alors réparties entre plusieurs locaux et mêlées à des dossiers hétérogènes. À partir de 1959, celles-ci sont reclassées selon un plan propre aux archives postérieures à la création de la Faculté impériale, élaboré par le doyen Giraud et organisé autour de quatre ensembles fonctionnels : administration générale, scolarité, personnel et comptabilité. Les documents préparatoires, la correspondance et le plan méthodique conservés permettent aujourd'hui de documenter cette opération\footcite{AHFM1MED208}\footcite{Perroud2009FondsModerne}. Ce plan est révisé en 1979.

Le traitement s'accompagne d'aménagements matériels. Les archives sont progressivement transférées dans la salle dite des \enquote{mâchicoulis} et dans les espaces attenants à l'actuelle salle Théophraste-Renaudot ; de nouveaux rayonnages sont installés et une poulie facilite, à partir de 1960, le déplacement des archives depuis les services. Ces transformations sont documentées par les dossiers de travaux et les photographies du fonds Giraud\footcite{BatisseursFaculte2021}. Elles montrent que le classement ne peut être dissocié des capacités de stockage et des circulations internes au bâtiment.

\subsection{La professionnalisation de la gestion des archives universitaires (1960--2022)}

Après le départ de Gaston Giraud, plusieurs agents poursuivent le classement, l'entretien et la communication des archives. Les annotations au crayon bleu et les perforations visibles sur les dossiers témoignent encore des traitements matériels de cette période. Aucun inventaire détaillé correspondant à la campagne initiale n'ayant été retrouvé, sa reconstitution repose sur le plan de classement, la correspondance administrative et l'examen des documents eux-mêmes.

En 1964, M\textsuperscript{lle} Roos, qui travaillait sous la direction d'Yvonne Vidal, est provisoirement remplacée par M. Tisseyre avant l'affectation d'un agent administratif\footcite{AHFM1MED208}. Le plan de classement des archives modernes et contemporaines élaboré par le doyen Gaston Giraud porte sur sa révision de 1979 les signatures d'André Coudurier, attaché d'administration responsable du service des \enquote{Affaires générales} et historien de l'enseignement médical, et de L.~Maury, responsable de la section \enquote{Archives}.\footnote{Leurs fonctions respectives sont indiquées sous leurs signatures dans le plan de classement mis à jour en 1979, en annexe de la feuille de route 2025-2029.} Ginette Ciurana assure ensuite la tenue des archives de 1989 à 2002, comme l'attestent les documents administratifs internes relatifs au fonctionnement de cette section\footcite{AHFMFeuilleRoute2025_2029}.

Ces interventions maintiennent l'accessibilité d'une partie des fonds, mais reposent sur des responsabilités individuelles et ne forment pas encore une politique archivistique à l'échelle de l'établissement. Cette situation est commune à de nombreuses universités françaises, longtemps dépourvues de services permanents malgré l'accroissement de leurs archives administratives et scientifiques\footcite{Maday2013}\footcite{Mercier2014}\footcite{GazetteArchivesESR2013}.

Une nouvelle étape s'ouvre au début des années 2000. Après un accompagnement des Archives départementales de l'Hérault pour les archives administratives de l'université et une visite de la directrice à l'UFR de médecine en 2004, l'Université recrute une archiviste en 2005, puis crée en 2008 un poste pérenne d'ingénieur d'études. Le service, successivement rattaché à la présidence puis à différentes directions, inscrit progressivement les archives dans l'organisation générale de l'établissement. \footcite{AHFMFeuilleRoute2025_2029} Cette professionnalisation accompagne un mouvement national de structuration des archives de l'enseignement supérieur et de mise en réseau de leurs responsables\footcite{Valencia2023}.

Un premier répertoire du fonds moderne est achevé en 2009. Il décrit alors 20,22 mètres linéaires d'archives relatives à l'administration générale entre 1794 et 1981 et identifie notamment, sous la cote 1~MED~208, les documents concernant la collecte et le classement des archives entre 1904 et 1968\footcite{Perroud2009FondsModerne}. Cet instrument confirme à la fois la continuité des interventions et les lacunes de l'historique de conservation.

Les archives historiques de médecine font ensuite l'objet d'un contrôle scientifique et technique en 2016, puis d'une mission de la Direction générale des patrimoines en 2019. Les échanges se poursuivent en 2022 sur leur maintien dans le bâtiment historique et sur les conditions nécessaires à leur conservation\footcite{AHFMReunionMars2022}\footcite{AHFMReunionJuillet2023}. Ces questions relèvent aujourd'hui notamment de la \gls{dcsph} (voir annexe~\ref{annexe:dcSPH}, fig.~\ref{fig:dcSPH-2025})\footcite{UMDCSPH2025}. Ces interventions n'aboutissent toutefois pas à une reprise générale du traitement avant la création de la Mission Archives historiques.

\subsection{La création de la Mission Archives historiques et la reprise du traitement (2023--)}

Une nouvelle phase commence le 12 avril 2023 lorsque Sophie Dikoff reçoit la responsabilité de la Mission Archives historiques de médecine. Cette formulation distingue la mission consacrée aux fonds de la Faculté du Service des archives de l'Université, dont elle assure par ailleurs la direction administrative. La reprise s'inscrit dans la réhabilitation du bâtiment ancien et dans une politique plus large de conservation et de valorisation du patrimoine documentaire universitaire\footcite{UMPatrimoineArchives}\footcite{UMServicesPatrimoine}\footcite{AHFMFeuilleRoute2025_2029}.

La création de cette mission rend possible la reprise d'un traitement d'ensemble des fonds ancien, moderne et contemporain. Elle ouvre une nouvelle phase, dont la formalisation juridique et opérationnelle relève de la préparation de la demande de dérogation.

\chapter{La conservation des archives historiques universitaires :
	cadre juridique et préparation de la dérogation au versement}
\chapitrecourt{Cadre juridique et dérogation au versement}
\label{chap:cadre-juridique-derogation}

Le maintien des Archives historiques dans le bâtiment de la Faculté, hérité de cette histoire, doit désormais être inscrit dans un cadre juridique. Les archives publiques destinées à être conservées sans limitation de durée ont en principe vocation à rejoindre le service public d'archives compétent ; leur conservation par l'Université suppose donc la préparation d'une demande de dérogation au versement.

La procédure préparée avec les Archives départementales de l'Hérault conduit l'Université à définir précisément le périmètre des fonds concernés et à réunir les garanties nécessaires à leur gestion, leur conservation et leur communication. Ce chapitre en présente le cadre juridique, les conditions de préparation et les conséquences en matière d'accès.

\section{Le périmètre et le cadre juridique de la dérogation}

\subsection{La coexistence d'archives publiques et de fonds privés}

Le Code du patrimoine définit les archives comme l'ensemble des documents et des données produits ou reçus dans l'exercice d'une activité, quels qu'en soient la date, le lieu de conservation, la forme et le support\footcite[art.~L.~211-1]{CodePatrimoine}. Leur statut dépend ainsi de leur provenance, non de leur ancienneté ni de leur lieu de conservation.

L'Université de Montpellier est un établissement public à caractère scientifique, culturel et professionnel\footcite[art.~L.~711-1]{CodeEducation}. Les documents produits ou reçus par la Faculté de médecine dans l'exercice de ses missions relèvent par conséquent des archives publiques\footcite[art.~L.~211-4]{CodePatrimoine}. Cette qualification concerne notamment les actes relatifs au gouvernement de l'institution, à la délivrance des grades, à la scolarité, aux enseignements, aux personnels, aux finances et aux relations avec les autorités universitaires ou administratives.

Les Archives historiques forment toutefois un ensemble plus composite. Aux archives institutionnelles se sont ajoutés des fonds de professeurs et de sociétés savantes, ainsi que des collections iconographiques constituées selon des modalités diverses. Leur présence dans les mêmes locaux ne leur confère pas un régime uniforme. Pour chaque ensemble doivent donc être établis le producteur, le mode d'entrée et, lorsqu'elles existent, les conditions particulières attachées au don, au legs ou au dépôt. Peuvent également intervenir des restrictions tenant à la vie privée, aux correspondances ou aux droits attachés aux œuvres reproduites\footcite{Foissey2017Juristes}.

La distinction des provenances permet d'appliquer à chaque fonds les règles qui lui correspondent et de renseigner son historique de conservation, ses modalités d'entrée, son statut ainsi que ses éventuelles restrictions de consultation, de reproduction ou de réutilisation. La demande de dérogation doit par conséquent porter sur un périmètre précisément identifié et distinguer les archives publiques auxquelles s'applique l'obligation de versement des fonds privés conservés à leurs côtés.

\subsection{Une dérogation à la destination des archives publiques}

Les archives publiques sélectionnées pour être gardées sans limitation de durée doivent en principe être versées dans un service public d'archives\footcite[art.~L.~212-4, I]{CodePatrimoine}. Pour l'Université de Montpellier, cette destination est normalement constituée par les Archives départementales de l'Hérault.

La sélection intervient à l'issue de la durée pendant laquelle les documents restent nécessaires aux services qui les ont produits. L'instruction interministérielle du 22 février 2005 fournit aux établissements relevant de l'Éducation nationale un cadre pour déterminer, selon les catégories documentaires, leur durée d'utilité administrative et leur sort final : conservation, élimination ou échantillonnage\footcite{InstructionArchivesEducation2005}. Dans le cas de la Faculté, son application doit tenir compte des transformations institutionnelles, de la continuité de certaines séries et de l'état inégal du classement des archives modernes et contemporaines.

L'article L.~212-4 permet toutefois à l'administration des archives d'en laisser la conservation au service compétent de l'établissement qui les a produites\footcite[art.~L.~212-4, I]{CodePatrimoine}. Cette possibilité est subordonnée à une convention précisant le périmètre concerné, les conditions de gestion, de conservation et de communication au public, les prescriptions scientifiques et techniques applicables ainsi que l'emploi d'une personne responsable qualifiée en archivistique\footcite[art.~R.~212-12]{CodePatrimoine}.

Ce dispositif est déjà appliqué à plusieurs institutions publiques. Les conventions publiées pour l'Académie des sciences, la Bibliothèque nationale de France, la Régie autonome des transports parisiens et l'Office français de protection des réfugiés et apatrides montrent la diversité des organismes susceptibles de conserver leurs archives définitives sous le contrôle de l'administration des archives.\footnote{Ces précédents ne sont pas entièrement assimilables au cas montpelliérain. L'Académie des sciences constitue le rapprochement institutionnel le plus pertinent en raison de son ancienneté et de la dimension scientifique et patrimoniale de ses archives \parencite{ConventionAcademieSciences2024}. La Bibliothèque nationale de France présente également un contexte patrimonial, mais sa convention concerne les archives produites par un établissement national \parencite{ConventionBnF2025}. Les conventions de la RATP et de l'OFPRA répondent à des situations administratives et documentaires différentes : la première porte sur les archives d'un opérateur de transport public, la seconde sur celles d'un établissement dont les dossiers sont étroitement liés à la protection des personnes \parencite{ConventionRATP2023,ConventionOFPRA2025}. Ces quatre conventions sont conclues avec le Service interministériel des Archives de France pour des documents qui auraient normalement vocation à rejoindre les Archives nationales. À l'échelle territoriale, la convention concernant le Pôle métropolitain d'Alsace relève d'une autre configuration : ses archives sont déposées auprès du service des Archives de l'Eurométropole de Strasbourg plutôt que maintenues directement par leur producteur ; le contrôle scientifique et technique demeure exercé par les Archives d'Alsace \parencite{ConventionPMAEurometropole2024}.}

La dérogation modifie ainsi la destination des documents, non leur statut. Les archives maintenues à la Faculté resteraient soumises aux règles applicables aux archives publiques ainsi qu'au contrôle scientifique et technique de l'État.

\subsection{Le contrôle scientifique et technique exercé par les Archives départementales}

Le contrôle scientifique et technique porte sur la gestion, la sélection, l'élimination, le classement, la conservation et la communication des archives publiques, qu'elles soient encore détenues par les services producteurs, maintenues par dérogation ou versées\footcite[art.~R.~212-2 et R.~212-3]{CodePatrimoine}. Dans le département, il est exercé par le directeur du service départemental d'archives et par les agents de l'État mis à disposition\footcite[art.~R.~212-4]{CodePatrimoine}.

Les Archives départementales de l'Hérault occupent donc une double position : elles constituent la destination réglementaire des archives publiques définitives de l'Université et apprécient les conditions proposées pour leur maintien à la Faculté. À ce titre, elles formulent des préconisations, examinent les instruments de recherche, autorisent les éliminations réglementaires et suivent l'avancement des opérations.

Ce suivi repose notamment sur la transmission de bilans périodiques et des instruments de recherche achevés\footcite{AHFMBilanSemestriel2025}\footcite{AHFMPointEtapeSeptembre2025}. Il permet d'adapter la programmation et d'évaluer si les moyens et les procédures mis en place peuvent être maintenus durablement, avant le dépôt de la demande.

\section{De l'expertise des fonds à la préparation de la demande}

\subsection{La formalisation d'une programmation pluriannuelle}

Les préconisations formulées au cours des interventions présentées au chapitre précédent sont réunies dans la feuille de route validée en octobre 2024 pour la période 2025--2029\footcite{AHFMFeuilleRoute2025_2029}. Celle-ci fixe les responsabilités, les moyens, les opérations à conduire et leurs modalités d'évaluation. Une programmation complémentaire, validée en avril 2025 puis actualisée en juillet, répartit les travaux à court, moyen et long termes\footcite{AHFMProgrammation2025}. La demande de dérogation constitue l'aboutissement envisagé de ce programme, et non une autorisation déjà acquise.

\subsection{Un traitement progressif adapté à l'état des fonds}

Les ensembles conservés présentent des états de traitement très différents. Le fonds ancien bénéficie d'un cadre de classement déjà établi, mais son récolement et la reprise de certaines séries restent nécessaires. Les archives modernes et contemporaines, les fonds privés et les collections iconographiques disposent quant à eux d'instruments de recherche de nature et de précision variables. Certains ensembles ont récemment été classés ou repris ; d'autres ne sont encore que partiellement décrits ou repérés\footcite{AHFMProgrammation2025}.

Les opérations procèdent donc par fonds ou sous-fonds cohérents, selon leur état, leur volume et les besoins de communication. Elles associent récolement, classement, reconditionnement et production d'instruments de recherche, sans attendre l'achèvement d'un classement général. Cette méthode permet de stabiliser progressivement les descriptions nécessaires à la gestion et à l'accès aux documents.

Une reprise particulière concerne la série dite \enquote{Q Supplément}, créée par Joseph Calmette pour accueillir, après son classement, les pièces antérieures à 1808 qui pourraient être retrouvées ou réintégrer ultérieurement les archives. Au fil du temps, elle a cependant reçu des documents de nature et de provenance diverses qui ne répondent pas tous à cette fonction initiale. Son récolement et la révision de son inventaire doivent donc permettre d'en préciser la composition et de rattacher, lorsque cela est possible, les documents aux fonds et aux instruments de recherche appropriés\footcite{AHFMQSupplement2025}.

Le traitement intellectuel est accompagné du reconditionnement progressif des ensembles, du suivi des conditions climatiques, de constats d'état, de restaurations et de la préparation d'un plan de sauvegarde. La réhabilitation du bâtiment doit parallèlement permettre d'aménager des espaces adaptés à la conservation, au travail des archivistes et à la consultation\footcite{AHFMFeuilleRoute2025_2029}.

\section{La communication des archives, condition de leur maintien}

\subsection{De la communicabilité juridique à l'accès effectif}

La conservation au sein de la Faculté ne modifie pas les règles de communicabilité et doit être distinguée de la dérogation individuelle permettant, sous certaines conditions, la consultation de documents avant l'expiration des délais légaux\footcite[art.~L.~213-1 à L.~213-3]{CodePatrimoine}. Les documents librement communicables doivent pouvoir être consultés auprès de l'Université ; les autres restent soumis aux restrictions et aux procédures prévues par le Code du patrimoine.

Les archives modernes et contemporaines peuvent contenir des informations relatives à la scolarité, aux carrières, à la vie privée ou à la santé. Leur classement doit permettre de repérer les catégories de données protégées au niveau pertinent. Une restriction appliquée indistinctement à un fonds entier pourrait rendre inaccessibles des documents librement communicables ; inversement, une description trop sommaire risquerait de laisser ignorer la présence d'informations soumises à des délais particuliers\footcite{Bastien1996DroitArchives}\footcite{Foissey2017Juristes}.

L'accès suppose également que le public puisse connaître la consistance des fonds et identifier les documents susceptibles de l'intéresser. Or le repérage repose encore largement sur la médiation des archivistes : les demandes sont fréquemment formulées à partir d'une personne, d'une date ou d'un sujet, sans référence préalable à une série ou à une cote (voir \cref{subsec:acces-documents-fonds}). La publication progressive des instruments stabilisés permettrait aux chercheurs d'effectuer eux-mêmes un premier repérage, sans supprimer l'accompagnement nécessaire aux recherches plus complexes.

\subsection{Une diffusion et une numérisation progressives}

Les instruments suffisamment stabilisés peuvent être publiés progressivement, à condition de préciser leur degré d'achèvement et les limites de l'offre disponible. La feuille de route prévoit également la mise en ligne de documents numérisés et le développement d'actions de valorisation\footcite{AHFMFeuilleRoute2025_2029}.

La numérisation doit être programmée selon les usages, l'état matériel et la communicabilité des documents. Elle peut faciliter l'accès aux pièces fréquemment demandées, limiter la manipulation des plus fragiles et soutenir des projets scientifiques ou patrimoniaux. Les registres régulièrement sollicités constituent à ce titre un corpus prioritaire possible\footcite{AHFMComitePilotage2026}.

La dispersion des sources relatives à l'enseignement médical entre la Faculté et les Archives départementales exige enfin d'articuler leur \gls{signalement}. Des renvois réciproques rendraient visibles des ensembles complémentaires sans remettre en cause leurs lieux de conservation respectifs\footcite{AHFMComitePilotage2026}.

\section*{Conclusion du chapitre et de la première partie}
\phantomsection
\label{sec:conclusion-partie-i}

Les Archives historiques résultent des transformations successives des institutions chargées de l'enseignement médical, mais aussi des usages, déplacements, pertes et reclassements qui ont affecté leurs documents. Leur organisation actuelle ne constitue donc pas un héritage demeuré intact : elle porte la trace des conditions dans lesquelles les fonds ont été produits, conservés et progressivement constitués comme patrimoine documentaire.

Leur maintien à la Faculté s'inscrit désormais dans un cadre juridique précis. La dérogation au versement ne peut reposer sur la seule ancienneté des fonds ni sur leur valeur patrimoniale ; elle suppose de garantir leur traitement, leur conservation, leur communicabilité et leur accessibilité sous le contrôle scientifique et technique de l'État. La préparation engagée jusqu'en 2029 conduit ainsi à articuler classement, récolement, amélioration des conditions matérielles, production des instruments de recherche et organisation de la consultation.

Rendre les fonds accessibles implique dès lors de dépasser leur seule conservation physique. Il faut permettre d'en connaître la composition, de repérer les documents pertinents et d'articuler les descriptions aux reproductions éventuellement disponibles, sans dissocier ces opérations du traitement archivistique. La deuxième partie examine l'environnement documentaire susceptible d'assurer cette continuité et d'organiser la diffusion des Archives historiques dans les infrastructures institutionnelles et les réseaux de \gls{signalement}.